\section{Algorithm}\label{sec:solution}

\ag{Explain the quality function and why it is a good proxy for the optimization function that you presented in \Cref{sec:model}, add proposition and proof for its sensitivity, add the algorithms (use the one-shot-top-k mechanism which is equivalent to $k$ times the exponential mechanism, see~\cite{TaoGMR22}), write the pseudo code, add a textual explanation for the code, and a proposition and proof for the algorithm being DP.}

\subsection{Overview}
\label{subsec:algorithm-overview}

The algorithm follows the model introduced in Section~\ref{sec:model}.
It first privately selects a small set of useful predicate views, then constructs explanation candidates by pairing selected views with attributes, and finally privately selects and releases the top-$k$ histogram explanations.
This structure mirrors recent differentially private explanation frameworks, which separate the definition of explanation quality from the design of low-sensitivity private selection mechanisms \cite{TaoGMR22}.
\alon{TODO: Add citation to Differentially Private Explanations for Clusters.}

Designing a practical algorithm for private query-difference explanations requires addressing two central challenges.
First, generating all possible views is infeasible because the space of predicate-defined subpopulations is combinatorial.
We therefore work with a candidate set $\mathcal{P}$ of predicate-based views, as is common in explanation and subgroup-discovery frameworks \cite{FEDEX,ICDE23}.
Second, the ideal utility function from Section~\ref{sec:explanation-quality} has high sensitivity and may be expensive to optimize.
We therefore introduce low-sensitivity proxy measures that capture the same conceptual goals: Histogram Distance for interestingness, Usefulness for representativeness and coverage, and DiversityScore for diversity.

\begin{figure}[ht]
\centering
\includegraphics[width=0.8\linewidth]{figures/placeholder.png}
\caption{An illustration of the private query-difference explanation framework.}
\label{fig:query-refinement}
\end{figure}
\alon{TODO: Replace the placeholder with the final framework figure. The figure should show: private database and queries $Q_1,Q_2$; private view selection using Usefulness; candidate construction $(p,A)$; private top-$k$ explanation selection using HD and DiversityScore; noisy histogram release.}

\subsection{Low-Sensitivity Proxy Measures}
\label{subsec:proxy-measures}

The proxy measures defined below are low-sensitivity counterparts of the ideal explanation-quality criteria introduced in Section~\ref{sec:explanation-quality}.
Each proxy is designed to capture one component of the model-level objective while remaining suitable for differentially private selection.
The View Usefulness score is used for private view selection; the Histogram Distance and DiversityScore are used for private explanation selection.

\subsubsection{View Usefulness}
\label{subsubsec:view-usefulness}

The ideal objective uses representativeness and coverage to prefer views that are both meaningful and statistically supported.
We capture these two criteria using a low-sensitivity View Usefulness score.
Coverage favors views that contain enough evidence from the compared query outputs:
\begin{equation}
Coverage(p,Q_1,Q_2,D)
=
\frac{|Q_1(p(D))|+|Q_2(p(D))|}{|D|}.
\label{eq:coverage-proxy}
\end{equation}

Representativeness favors views whose attribute distributions are not too far from the distribution of the full database.
Let $P_D$ and $P_{p(D)}$ denote the empirical distributions used to compare the database and the view.
We define
\begin{equation}
Representativeness(p,D)
=
1-TVD(P_D,P_{p(D)}).
\label{eq:representativeness-proxy}
\end{equation}
The Usefulness score is then
\begin{equation}
Usefulness(p,Q_1,Q_2,D)
=
Representativeness(p,D) \cdot Coverage(p,Q_1,Q_2,D).
\label{eq:usefulness-proxy}
\end{equation}
\alon{TODO: Specify precisely which empirical distributions are compared in $Representativeness(p,D)$: all attributes, the candidate attribute $A$, or a fixed set of attributes. The previous draft used $TVD(D,p(D))$ directly; since higher utility should be better, the current version uses $1-TVD$. Verify this convention against the implementation.}

\begin{proposition}[Sensitivity of View Usefulness]
\label{prop:usefulness-sensitivity}
The View Usefulness score has bounded global sensitivity.
\end{proposition}

\begin{proof}
\alon{TODO: Complete proof.}
\end{proof}

\subsubsection{Histogram Distance}
\label{subsubsec:histogram-distance}

Histogram Distance serves as a proxy for interestingness.
In the ideal objective, an explanation is interesting when the distribution of an attribute changes substantially between $Q_1(p(D))$ and $Q_2(p(D))$.
Histogram Distance captures this idea by comparing the one-way histograms of the same attribute over the two query outputs restricted to the same predicate view.

\begin{definition}[Histogram Distance]
Given a dataset $D$, a view predicate $p$, an attribute $A$, and two queries $Q_1,Q_2$, let
\[
\operatorname{Dom}(p(D)[A])=\{v \in \operatorname{Dom}(D[A]) \mid H_{p(D),A}(v)>0\}
\]
be the set of non-empty histogram bins for $A$ inside $p(D)$.
The \emph{Histogram Distance} is defined as
\begin{equation}
HD(D;p,A,Q_1,Q_2)
=
\sum_{v \in \operatorname{Dom}(p(D)[A])}
\frac{\left|H_{Q_1(p(D)),A}(v)-H_{Q_2(p(D)),A}(v)\right|}{H_{p(D),A}(v)}.
\label{eq:hd-definition}
\end{equation}
\end{definition}

The denominator normalizes the contribution of each bin by its size within the view, so the score emphasizes relative changes rather than only absolute count differences.
Empty bins are excluded to avoid division by zero.
\alon{TODO: Verify that excluding empty bins is the desired convention. Alternatively, define smoothing with $H_{p(D),A}(v)+\gamma$. This choice affects the sensitivity proof.}

\begin{proposition}[Sensitivity of Histogram Distance]
\label{prop:hd-sensitivity}
The Histogram Distance score has bounded global sensitivity.
\end{proposition}

\begin{proof}
Let $D, D'$ be neighboring datasets such that $|\,|D'| - |D|\,| = 1$ and $D' = D \cup \{t\}$,  
where $t$ is a single tuple.  
Let $t_A = V'$ denote the bucket value for attribute $A$ of tuple $t$.

The global sensitivity of $HD$ is:
\begin{equation}
GS_{HD} = \max_{\substack{D,D' \\ D \sim D'}} 
\left| HD(D) - HD(D') \right|,
\end{equation}
where $D \sim D'$ denotes that $D$ and $D'$ are neighboring datasets.

From~\eqref{eq:hd-definition}, we have:
\begin{align}
&\left| HD(D) - HD(D') \right| \notag \\
&= \bigg| 
\sum_{V \in \text{buckets}(D'[A])} 
\frac{\left| H_{Q(p(D)),A}(V) - H_{Q'(p(D)),A}(V) \right|}
     {H_{p(D),A}(V)} \notag \\
&\quad - 
\sum_{V \in \text{buckets}(D'[A])} 
\frac{\left| H_{Q(p(D')),A}(V) - H_{Q'(p(D')),A}(V) \right|}
     {H_{p(D'),A}(V)}
\bigg|.
\end{align}

We partition the set of buckets into two disjoint subsets:
\[
\text{buckets}(D'[A]) = 
\left(\text{buckets}(D'[A]) \setminus \{V'\}\right) \cup \{V'\}.
\]

\textbf{Case 1: $V \neq V'$}.  
For all $V \neq V'$, the numerator in each term is unchanged between $D$ and $D'$,  
while the denominator changes by at most $1$ because the tuple $t$ is not assigned to $V$.  
Thus:
\[
\sum_{V\ne V' \in \text{buckets}(D'[A])} \left| 
\frac{\cdot}{H_{p(D),A}(V)} - \frac{\cdot}{H_{p(D'),A}(V)}
\right| \leq 1.
\]

\textbf{Case 2: $V = V'$}.  
Here, both numerator and denominator may change, but by the neighboring dataset property,  
each changes by at most $1$. Therefore:
\begin{align}
&\left| 
\frac{\left| H_{Q(p(D)),A}(V') - H_{Q'(p(D)),A}(V') \right|}{H_{p(D),A}(V')}
\right. \notag \\
&\quad \left.
-
\frac{\left| H_{Q(p(D')),A}(V') - H_{Q'(p(D')),A}(V') \right|}{H_{p(D'),A}(V')}
\right| \leq 1.
\end{align}



Since only one bucket $V'$ can change between $D$ and $D'$,  
and the change in that bucket’s contribution is bounded by $1$,  
it follows that:
\[
\left| HD(D) - HD(D') \right| \leq 1.
\]

Taking the maximum over all neighboring datasets yields:
\[
GS_{HD} \leq 1.
\]
\end{proof}

\subsubsection{Diversity Score}
\label{subsubsec:diversity-score}

DiversityScore corresponds to the diversity criterion from the model section.
When multiple explanations are released, DiversityScore discourages repeatedly selecting explanations that focus on the same attribute or highly similar aspects of the query difference.

For a candidate $(p,A)$ and a set $S$ of attributes selected in previous iterations, we use the following simple penalty:
\begin{equation}
DiversityScore(A,S)
=
\begin{cases}
1, & A \notin S,\\
\lambda, & A \in S,
\end{cases}
\label{eq:diversity-score}
\end{equation}
where $0\leq \lambda \leq 1$ controls the penalty for repeatedly selecting the same attribute.
Setting $\lambda=0$ prevents the same attribute from being selected more than once, while larger values allow repeated attributes with a penalty.
Since $DiversityScore$ depends only on previously selected outputs and not directly on the private database, it does not increase the data sensitivity of the selection utility.
\alon{TODO: Decide whether diversity should be attribute-level only, view-level, or based on a distance between histograms. If a richer diversity definition is used, update the sensitivity argument accordingly.}

\subsubsection{Combined Proxy Objective}
\label{subsubsec:combined-proxy-objective}

The proxy measures are used in two stages.
First, View Usefulness accounts for representativeness and coverage during private view selection.
Second, for private explanation selection over candidates $c=(p,A)$, we use
\begin{equation}
\widetilde{U}_i(D;c)
=
HD(D;p,A,Q_1,Q_2) \cdot DiversityScore(A,S_{i-1}),
\label{eq:selection-utility}
\end{equation}
where $S_{i-1}$ is the set of attributes selected in the previous iterations.
Thus, the final selection utility uses Histogram Distance as a proxy for interestingness and DiversityScore as a proxy for diversity, while representativeness and coverage have already been handled by the view-selection stage.

\begin{proposition}[Sensitivity of the selection utility]
\label{prop:selection-utility-sensitivity}
The selection utility $\widetilde{U}_i$ has bounded global sensitivity whenever $HD$ has bounded global sensitivity.
\end{proposition}

\begin{proof}
The diversity term is bounded in $[0,1]$ and, conditioned on the previous selected outputs, does not directly access the private database.
Therefore it does not increase the global sensitivity of the data-dependent term $HD$.
\alon{TODO: Complete the formal argument under adaptive composition. The proof should explicitly condition on the history of previously selected candidates.}
\end{proof}

\subsection{Hardness of Non-Private Proxy Explanation Selection}
\label{subsec:non-private-hardness}

Even without privacy, selecting the best set of explanations under the proxy objective can be computationally difficult because the algorithm must choose a subset of candidates that balances quality and diversity over a combinatorial candidate space.

\begin{theorem}[Hardness of non-private proxy explanation selection]
\label{thm:non-private-hardness}
The problem of selecting the optimal set of $k$ explanation candidates under the proxy objective is NP-hard, even when the proxy scores are known exactly and no privacy constraint is imposed.
\end{theorem}

\begin{proof}
\alon{TODO: Add reduction, likely from Maximum Coverage or Maximum Diversity.}
\end{proof}

\subsection{Private Query-Difference Explanation Algorithm}
\label{subsec:private-algorithm}

Algorithm~\ref{alg:main} gives the end-to-end pipeline.
It receives the private database, the two related queries, a candidate set of predicate views, the attribute set, the number of views $\tau$, the number of explanations $k$, and the privacy budget $\epsilon$.
It allocates the privacy budget across private view selection, private explanation selection, and noisy histogram release.

\begin{algorithm}[t]
\caption{Private-Queries-Diff-Top-$k$-Explanation}
\label{alg:main}
\begin{algorithmic}[1]
\Require Private dataset $D$; queries $Q_1, Q_2$; predicate-based views $\mathcal{P}$; attribute set $Attrs$; number of selected views $\tau$; number of explanations $k$; privacy budget $\epsilon$
\Ensure A set of $k$ differentially private explanation histograms

\Statex \textbf{Overview.} Privately select a subset of informative views, generate candidates, and privately select and release the top-$k$ explanations.
\State \textbf{/* Privacy budget allocation */}
\State $\epsilon_P, \epsilon_S, \epsilon_H  \gets \epsilon/2, \epsilon/4, \epsilon/4$

\State \textbf{/* Private view selection */}
\State Define $u_{use}(p) \gets Usefulness(p, D, Q_1, Q_2)$
\State $\widehat{\mathcal{P}} \gets \textsc{OneShotTop-}\tau(\mathcal{P}, u_{use}(\cdot), \epsilon_P)$ \Comment{$\epsilon_P$-DP}

\Statex \textbf{(Post-processing)} Candidate construction does not incur additional privacy loss.

\State \textbf{/* Construct explanation candidates */}
\State $C \gets \emptyset$
\ForAll{$p \in \widehat{\mathcal{P}}$}
    \ForAll{$A \in Attrs$}
        \State $C \gets C \cup \{(p,A)\}$
    \EndFor
\EndFor

\State \textbf{/* Private top-$k$ explanation selection and release */}
\State $\mathcal{E} \gets \textsc{Find-Top-$k$-Explanations}(D,Q_1,Q_2,C,k,\epsilon_S,\epsilon_H)$

\Return $\mathcal{E}$
\end{algorithmic}
\end{algorithm}

The algorithm has three stages.
First, it privately selects useful views using the Usefulness proxy.
Second, it constructs candidate explanations by pairing each selected view with each attribute; this is post-processing.
Third, it invokes Algorithm~\ref{alg:topk} to privately select and release the final top-$k$ explanations.

\begin{algorithm}[t]
\caption{Find-Top-$k$-Explanations}
\label{alg:topk}
\begin{algorithmic}[1]
\Require Private dataset $D$; queries $Q_1,Q_2$; candidates $C$; number of explanations $k$; budgets $\epsilon_S,\epsilon_H$
\Ensure $k$ differentially private explanation histograms

\State $ChosenAttrs \gets \emptyset$
\State $\mathcal{E} \gets \emptyset$

\For{$i \gets 1$ to $k$}
  \State \textbf{/* Private selection using the exponential mechanism */}
  \State $u_i(c=(p,A)) \gets HD(D;p,A,Q_1,Q_2) \cdot DiversityScore(A,ChosenAttrs)$
  \State $c_i=(p_i,A_i) \gets \textsc{ExponentialMechanism}(C,u_i(\cdot),\epsilon_S/k)$

  \State \textbf{/* Private histogram release */}
  \State $\Delta H_i(v) \gets H_{Q_1(p_i(D)),A_i}(v)-H_{Q_2(p_i(D)),A_i}(v)$ for all bins $v$
  \State $\widetilde{\Delta H_i}(v) \gets \Delta H_i(v)+\textsc{Laplace}(2k/\epsilon_H)$ for all bins $v$

  \State $\mathcal{E} \gets \mathcal{E} \cup \{(p_i,A_i,\widetilde{\Delta H_i})\}$
  \State $ChosenAttrs \gets ChosenAttrs \cup \{A_i\}$
  \State $C \gets C \setminus \{c_i\}$
\EndFor

\State \Return $\mathcal{E}$
\end{algorithmic}
\end{algorithm}

Algorithm~\ref{alg:topk} proceeds iteratively for $k$ rounds.
In each iteration, it computes a selection utility for every remaining candidate using Histogram Distance and the diversity penalty, selects one candidate privately, and releases a noisy difference histogram for the selected candidate.
The difference histogram captures how the distribution of attribute $A_i$ changes between $Q_1(p_i(D))$ and $Q_2(p_i(D))$.

\subsection{Differential Privacy Analysis}
\label{subsec:dp-proof}

We now prove that Algorithms~\ref{alg:main} and~\ref{alg:topk} satisfy $\epsilon$-differential privacy with respect to the input private dataset $D$.
The proof uses the sensitivity bounds for the proxy scores, the exponential mechanism for private selection, the Laplace mechanism for histogram release, and sequential composition.

\begin{proposition}[Sensitivity of noisy difference histogram release]
\label{prop:histogram-release-sensitivity}
For a fixed candidate $(p,A)$, the difference histogram
\[
\Delta H_{p,A}(v)=H_{Q_1(p(D)),A}(v)-H_{Q_2(p(D)),A}(v)
\]
has bounded $L_1$-sensitivity.
\end{proposition}

\begin{proof}
\alon{TODO: Complete proof.}
\end{proof}

\begin{proposition}[Differential Privacy of Algorithm~\ref{alg:main}]
\label{thm:alg1-dp}
Algorithm~\ref{alg:main} satisfies $\epsilon$-differential privacy.
\end{proposition}

\begin{proof}
Algorithm~\ref{alg:main} partitions the privacy budget into three components: $\epsilon_P$ for private view selection, $\epsilon_S$ for private explanation selection, and $\epsilon_H$ for private histogram release, with $\epsilon_P+\epsilon_S+\epsilon_H=\epsilon$.

The first privacy-sensitive step is the selection of predicate-based views using a One-Shot Top-$\tau$ mechanism with the Usefulness score.
By Proposition~\ref{prop:usefulness-sensitivity}, this score has bounded sensitivity, so the view-selection step satisfies $\epsilon_P$-differential privacy.
\alon{TODO: Cite the exact one-shot top-$k$ or one-shot top-$\tau$ mechanism used here and state its sensitivity parameter.}

Candidate construction pairs selected views with attributes.
This step is a deterministic function of the privately selected views and therefore incurs no additional privacy loss by post-processing.
Algorithm~\ref{alg:main} then invokes Algorithm~\ref{alg:topk}, which satisfies $(\epsilon_S+\epsilon_H)$-differential privacy by Proposition~\ref{thm:alg2-dp}.
By sequential composition, the total privacy loss is $\epsilon_P+\epsilon_S+\epsilon_H=\epsilon$.
\end{proof}

\begin{proposition}[Differential Privacy of Algorithm~\ref{alg:topk}]
\label{thm:alg2-dp}
Algorithm~\ref{alg:topk} satisfies $(\epsilon_S+\epsilon_H)$-differential privacy.
\end{proposition}

\begin{proof}
Algorithm~\ref{alg:topk} performs $k$ adaptive iterations.
In each iteration, it first selects an explanation candidate using the exponential mechanism with budget $\epsilon_S/k$.
The utility used by the exponential mechanism is $\widetilde{U}_i$, whose sensitivity is bounded by Proposition~\ref{prop:selection-utility-sensitivity}.
Therefore, each selection step satisfies $(\epsilon_S/k)$-differential privacy.

After selecting a candidate, the algorithm releases a noisy difference histogram using the Laplace mechanism with budget $\epsilon_H/k$.
By Proposition~\ref{prop:histogram-release-sensitivity}, the released histogram has bounded $L_1$-sensitivity, so each histogram release satisfies $(\epsilon_H/k)$-differential privacy.

By sequential composition within each iteration and across the $k$ iterations, the total privacy loss is
\[
    k\left(\frac{\epsilon_S}{k}+\frac{\epsilon_H}{k}\right)
    =
    \epsilon_S+\epsilon_H.
\]
The iterations are adaptive because later utilities depend on previously selected attributes, but adaptive sequential composition still applies.
All remaining operations, including updating the candidate set and the selected-attribute set, are post-processing of previously released private outputs.
Thus, Algorithm~\ref{alg:topk} satisfies $(\epsilon_S+\epsilon_H)$-differential privacy.
\alon{TODO: Add citation to adaptive composition, e.g., \cite{DBLP:series/synthesis/2016Li}, and verify whether the noisy histogram release should be analyzed under parallel or sequential composition if histograms use disjoint views.}
\end{proof}

\begin{corollary}[End-to-End Differential Privacy]
\label{cor:end-to-end-dp}
By sequential composition of Algorithms~\ref{alg:main} and~\ref{alg:topk}, the overall explanation pipeline satisfies $\epsilon$-differential privacy.
\end{corollary}

\subsection{Complexity Analysis}
\label{subsec:complexity-analysis}

Let $|\mathcal{P}|$ be the number of candidate predicate views, $m=|Attr(D)|$ the number of attributes, and $b$ the maximum number of bins in an attribute histogram.
The view-selection stage evaluates Usefulness over the candidate views.
After selecting $\tau$ views, the algorithm constructs at most $\tau m$ explanation candidates.
Each iteration of Algorithm~\ref{alg:topk} evaluates the selection utility over the remaining candidates and computes one noisy histogram.
\alon{TODO: Complete the asymptotic complexity bound in terms of $|D|$, $|\mathcal{P}|$, $m$, $b$, $\tau$, and $k$. Also explain whether histograms are precomputed or computed on demand.}
